\documentclass[11pt]{amsart}
\usepackage[margin=1.1in]{geometry}
\usepackage{amsmath,amssymb,amsthm,mathtools}
\usepackage{enumitem}
\usepackage{microtype}
\usepackage{hyperref}
\usepackage[nameinlink,capitalize]{cleveref}
\hypersetup{colorlinks=true,linkcolor=blue,citecolor=blue,urlcolor=blue}

\theoremstyle{plain}
\newtheorem{theorem}{Theorem}[section]
\newtheorem{lemma}[theorem]{Lemma}
\newtheorem{proposition}[theorem]{Proposition}
\theoremstyle{definition}
\newtheorem{definition}[theorem]{Definition}
\theoremstyle{remark}
\newtheorem{remark}[theorem]{Remark}

\title[Local Repaired-Gauge Representation]
{Local Repaired-Gauge Representation\\
for Single-Core Type II Navier--Stokes Concentration}
\author{}
\date{}

\begin{document}
\begin{abstract}
We prove the local representation theorem needed to place a genuine single-core
Type II Navier--Stokes concentration into the associated compactness criterion.  Starting from a positive local CKN core with compact-cylinder suitable
bounds and a unique nondegenerate scale and center, we construct a local
renormalized chart, compactly supported repaired gauges, local pressure
representation, and modulation coefficients.  The resulting repaired orbit
satisfies the local single-core compactness criterion.
The pressure reconstruction and the modulation coefficients are consequences
of the local chart and repaired gauge; they are not additional assumptions.
Failure of the localized gauge Jacobian or of compact-cylinder uniqueness
places the sequence in the multibubble or cascade configurations.
This paper uses the classical Navier--Stokes theory for suitable weak solutions (Leray \cite{Leray1934}, Caffarelli--Kohn--Nirenberg \cite{CaffarelliKohnNirenberg1982}, Scheffer \cite{Scheffer1976}, and Lin \cite{Lin1998}).
\end{abstract}
\maketitle

\section{Single-core input}

Let \((u,p)\) be a suitable weak solution near a possible singular point
\((x_0,T)\).  A local Type II concentration sequence is a sequence
\(r_n\downarrow0\) such that the parabolic rescalings
\[
u_n(y,s)=r_n u(x_0+r_ny,T+r_n^2s),
\qquad
p_n(y,s)=r_n^2p(x_0+r_ny,T+r_n^2s)
\]
are suitable on every fixed backward cylinder and retain positive local CKN
mass on one fixed cylinder.

For a suitable pair \((v,q)\), let \(A_v,E_v,C_{v,q},D_{v,q}\) denote the usual
scale-invariant local energy, dissipation, velocity, and pressure quantities on
backward cylinders.  The single-core case is the case in which:
\begin{enumerate}[label=\textup{(\roman*)}]
\item \(A+E+C+D\) is bounded on every compact rescaled cylinder;
\item one selected core carries a fixed positive CKN density;
\item outside a larger compact cylinder, no comparable active core remains in
      the local rescaled frame;
\item the localized scale and center selected by that core are unique in the
      sense that the localized scale-center Jacobian is invertible.
\end{enumerate}
If any of these conditions fails, the case is not single-core and belongs to
the multibubble/cascade analysis.

\section{Local chart and equation}

Choose absolutely continuous local scale and center functions \(\Lambda(\tau)>0\)
and \(X(\tau)\), and set
\[
Y=\frac{x-X(\tau)}{\Lambda(\tau)},
\qquad
\frac{d\tau}{dt}=\Lambda(\tau)^{-2},
\qquad
u(x,t)=\Lambda(\tau)^{-1}V(Y,\tau).
\]
On compact cylinders, the transformed velocity satisfies
\[
\partial_\tau V+(V\cdot\nabla)V+\nabla P
=\nu\Delta V+a(\tau)(V+Y\cdot\nabla V)+b(\tau)\cdot\nabla V,
\qquad \nabla\cdot V=0,
\]
where
\[
a(\tau)=\partial_\tau\log\Lambda(\tau),
\qquad
b(\tau)=-\Lambda(\tau)^{-1}\partial_\tau X(\tau)
\]
up to the harmless sign convention determined by the chart.  Only consistency
with the displayed equation is used below.

\begin{lemma}[Local change of variables]
\label{lem:chart}
Let \(\Lambda>0\) and \(X\) be absolutely continuous on a compact time
window, and define \(V,P,a,b\) by the chart above.  Then the transformed fields
solve the displayed renormalized Navier--Stokes equation in the sense of
distributions on compact cylinders.
\end{lemma}

\begin{proof}
This is the standard distributional change of variables for Navier--Stokes
under an absolutely continuous translation and dilation.  The terms containing
\(\partial_\tau\log\Lambda\) and \(\partial_\tau X\) are exactly the scale and
translation modulation terms in the displayed equation.
\end{proof}

\section{Localized repaired gauges}

Fix \(R\) and \(\chi_R\in C_c^\infty(B_{2R})\), \(\chi_R=1\) on \(B_R\).  Set
\[
G_{0,R}(V)=\int \chi_R(Y)|Y|^{-p}|V(Y)|^3\,dY-\Theta_0,
\qquad 0<p<3,
\]
\[
G_{j,R}(V)=\int Y_j\chi_R(Y)|V(Y)|^2\,dY,
\qquad j=1,2,3.
\]

\begin{lemma}[Scale row]\label{lem:scale-row}
For variations generated by scaling,
\[
DG_{0,R}(V)[V+Y\cdot\nabla V]
=p\int \chi_R|Y|^{-p}|V|^3
 -\int (Y\cdot\nabla\chi_R)|Y|^{-p}|V|^3 .
\]
Thus, if the cutoff-annulus term is strictly smaller than the weighted core
contribution, the scale row is transverse.
\end{lemma}

\begin{proof}
Differentiate \(G_{0,R}(\lambda V(\lambda\cdot))\) at \(\lambda=1\).  The
homogeneity of \(|Y|^{-p}\) gives the first term, and differentiating the cutoff
gives the annular error.
\end{proof}

\begin{theorem}[Local repaired-gauge solve]\label{thm:gauge-solve}
On every compact time window where the localized scale-center Jacobian is invertible with bounded inverse, there are absolutely continuous local scale and center corrections imposing the four constraints.  If this localized Jacobian is not invertible, the branch is not a nondegenerate single-core branch and belongs to the multibubble, cascade, or gauge-degenerate alternative.
\end{theorem}

\begin{proof}
The constraints are continuously differentiable with respect to the finite
number of local scale and center parameters on the compact support of \(\chi_R\).
The quantitative inverse function theorem gives a unique correction while the
Jacobian remains invertible.  Absolute continuity follows from the absolute
continuity of the unrepaired chart and the local integrability of the time
derivatives of the constraints.  Failure of the Jacobian means that the selected
mass does not determine one local scale and center, which is precisely the
non-single-core alternative.
\end{proof}

\section{Pressure and modulation}

For every compact ball, decompose
\[
P=P_{\rm loc}+H,
\qquad
-\Delta P_{\rm loc}=\partial_i\partial_j(\zeta V_iV_j),
\]
where \(\zeta\in C_c^\infty(B_{2R})\) equals one on the smaller ball.  Then
\(H\) is harmonic on the smaller ball.  Calderon--Zygmund estimates for
\(P_{\rm loc}\) and interior harmonic estimates for \(H\) give the pressure
bounds used in local CKN and Caccioppoli estimates.

\begin{proposition}[Modulation coefficients]\label{prop:modulation}
On every compact selected window on which the localized gauge Jacobian remains
invertible and the local Caccioppoli bounds hold, the repaired chart has
measurable modulation coefficients \(a,b\) that are locally bounded in terms of
the inverse Jacobian, the local pressure bounds, and the compact-cylinder
suitable norms.
\end{proposition}

\begin{proof}
Differentiate the four gauge constraints in time.  The coefficients \(a,b\)
appear linearly through the scale and translation rows of the renormalized
equation.  The coefficient matrix is the localized gauge Jacobian.  The forcing
terms consist of local convection, viscosity, and pressure contributions, all
estimated on the compact support of \(\chi_R\) by Caccioppoli and the local
pressure decomposition.  Inverting the Jacobian gives the asserted local bounds.
\end{proof}

\section{Representation theorem}

\begin{theorem}[Single-core repaired-gauge representation]\label{thm:representation}
Let a local Type II concentration sequence be single-core in the sense of
Section 1.  On every compact window on which the selected core determines an absolutely continuous local rescaled frame and the localized scale-center Jacobian is invertible, after passing to a subsequence the branch
admits a localized repaired-gauge orbit \((V,P,a,b)\) satisfying:
\begin{enumerate}[label=\textup{(\roman*)}]
\item the renormalized Navier--Stokes equation on compact cylinders;
\item positive selected local CKN mass on the selected core;
\item local pressure reconstruction modulo spatial means;
\item locally bounded modulation coefficients on the selected windows;
\item the compact-cylinder suitable and Caccioppoli bounds used by the
      local single-core compactness criterion.
\end{enumerate}
The pressure reconstruction and the coefficients \(a,b\) in items
\textup{(iii)}--\textup{(iv)} are conclusions of the theorem.  The only
structural alternatives to the conclusion are absence of a single absolutely
continuous local frame, failure of localized gauge nondegeneracy, or loss of
compact-cylinder suitable bounds; these are precisely the multibubble, cascade,
gauge-degenerate, or local compactness-failure cases.
\end{theorem}

\begin{proof}
The local rescaled frame and \cref{lem:chart} give the renormalized equation
before repair.  The repaired-gauge solve imposes the localized scale and
centering constraints and preserves absolute continuity of the corrected
parameters.  The pressure reconstruction follows from the Navier--Stokes
pressure equation under the same change of variables, modulo functions of time,
and the compact-ball decomposition gives the local pressure bounds.  The
coefficients \(a,b\) are forced by the absolutely continuous scale and
translation parameters in the transformed equation; \cref{prop:modulation}
gives their local bounds on selected windows.  The selected CKN mass and compact-cylinder suitable bounds define the nondegenerate compact single-core branch.  Thus the resulting orbit satisfies the local bounds used by the local single-core compactness criterion.  Failure of the localized scale/center uniqueness is the
non-single-core alternative.
\end{proof}

\bibliographystyle{alpha}
\bibliography{paper2_renormalized_typeII_repaired_gauges}

\end{document}
